\section{标注数据与树库（Treebanks）}

\subsection{从规则到数据驱动}

早期的 NLP 系统依赖\textbf{手写规则}来解析句子结构。这种方法存在两个致命问题：

\begin{enumerate}
    \item \textbf{覆盖率不足}：人类语言具有无穷的创造性（如 Yoda 式语序重排、网络用语等），手写规则难以穷举
    \item \textbf{无法消歧}：纯规则系统可以生成所有合法的句法分析结果，但无法在多个合法分析之间选择最可能的那个
\end{enumerate}

\begin{figure}[H]
\centering
\includegraphics[width=0.95\textwidth]{slides-images/slide-023.jpg}
\caption{从规则到数据驱动的转变\protect\footnotemark}
\end{figure}
\footnotetext{Slides 第23页。}

为了解决这些问题，NLP 社区从 1980 年代末到 1990 年代开始大规模建设\textbf{标注语料库}（annotated corpora），即\textbf{树库}（Treebanks）——由人工标注了句法结构的句子集合。

\begin{knowledgebox}{树库为何重要}
树库的建设虽然缓慢而费力，但带来了两大革命性优势：
\begin{enumerate}
    \item \textbf{可复用的训练数据}：一次标注，多次使用。可以训练句法分析器、词性标注器、心理语言学模型等
    \item \textbf{系统化的评估方法}：提供了客观的基准测试。在树库出现之前（1950s--1980s），NLP 系统的评估方式是：运行程序，输入一句话，看看输出的分析结果是否``看起来合理''——没有任何系统化的定量评估
\end{enumerate}
\end{knowledgebox}

\subsection{Universal Dependencies（UD）}

\begin{figure}[H]
\centering
\includegraphics[width=0.95\textwidth]{slides-images/slide-024.jpg}
\caption{Universal Dependencies 树库项目\protect\footnotemark}
\end{figure}
\footnotetext{Slides 第24页。}

Manning 介绍了他深度参与的 \textbf{Universal Dependencies}（UD）项目——一个统一标注方案的跨语言依存树库：

\begin{itemize}
    \item 覆盖超过 \textbf{100 种语言}
    \item 使用统一的依存关系标签体系
    \item 适用于跨语言研究、心理语言学实验等
    \item 课程作业 Assignment 2 使用的就是这类数据
\end{itemize}

\subsection{本章小结}

树库的建设推动了 NLP 从规则驱动转向数据驱动。Universal Dependencies 项目提供了覆盖 100 多种语言的标准化依存树库，是现代依存分析研究和应用的基础资源。树库同时提供了训练数据和评估基准。

%% ========== 基于转移的依存分析 ==========

